\documentclass{jsarticle}
\usepackage{amsmath}
\usepackage{amssymb}
\begin{document}
\title{行動による人が持っている情報ネットワークの \\
可能性の発見と再構築、情報空間 \\
によるモデル}
\author{苫米地博士の論文}
\date{}
\maketitle

      初速度のドーパミンの量からの行動の補助から、終わりの精神のドーパミン
  の量の調節、相加相乗平均から、始まりと終わりの精神のドーパミンの
  調節による人の行動の領域による予知
  $$ x \log x = \left(v_0 + {1 \over at_n}\right)V, V_x = e^f + {1 \over e^f} $$
     経路積分によるエントロピー不変量から
  このドーパミンの量から人の行動傾向が決まる。
  基礎代謝の持続できる距離と精神の介入と補助、言語の発生と抑制による経路
  パターンの組み合わせを再構築する。
  $$ ||dx^2|| = \int {\mathrm{exp}\left({1 \over \sqrt{2 \tau q}}L(x)
  \right) + \mathrm{O}(N^{-1})}d \tau   $$
  $$ p = {v \over 2m}, \hbar = {p \over 2mv}, \lambda = h \nu $$
     これらの式から、言語の神経ネットワーク
  による単体クラスの組み合わせパターンとその対象の人の領域による
  人が持っている情報空間の傾向の可能性のポテンシャルエネルギーの
  予知を人工知能の正規表現で抽出する。
   
  $$ {d \over d \gamma}\Gamma = \Gamma^{{\gamma}^{'}} = e^{-f}  $$
  $$ C = \int \left(\int {1 \over x^s}dx + \log x \right)\mathrm{dvol} $$
  オイラーの定数は、ゼータ関数とゼータ関数自体のエントロピーを
  体積積分した結果でもある。このオイラーの定数は、ガンマ関数の
  大域的微分方程式からも導出できる。このガンマ関数は
  大脳基底核のエントロピー不変量の強度も表している。
  このガンマ関数の大域的微分多様体による、ブレインインターフェース
  での人体のデータ計測もこのJones多項式とラプラス方程式からの
  ３次元多様体のエントロピー不変式からも、言語とイメージの
  大脳基底核からのデータ抽出も可能と言える。

\end{document}
